% !TEX root = ../thesis.tex

\chapter{Úvod}\label{ch:introduction}

Úvod práce stručne opisuje stanovený problém, kontext problému a motiváciu pre riešenie problému. Z úvodu by malo byť jasné, že stanovený problém doposiaľ nie je vyriešený a má zmysel ho riešiť.

V úvode neuvádzajte štruktúru práce, t.j. o čom je ktorá kapitola. Rozsah úvodu je minimálne 2 celé strany (vrátane formulácie úlohy).

\section{Motivácia}

V oblasti energetiky v Európe vzniká veľké množstvo odborných informácií, no ich praktické využitie je často pomalé a náročné na orientáciu. Pri riešení konkrétnej otázky musí používateľ prechádzať viacero zdrojov, porovnávať údaje a ručne vytvárať závery, čo predlžuje rozhodovanie a zvyšuje riziko nepresnej interpretácie. Tento problém je výrazný najmä v situáciách, keď je potrebné reagovať rýchlo a s obmedzenými personálnymi alebo technickými kapacitami.

Motiváciou práce je preto overiť, či je možné pomocou relatívne malých výpočtových zdrojov doučiť malé jazykové modely na úzko špecializovaných energetických dátach tak, aby poskytovali rýchle, vecné a prakticky použiteľné odpovede. Súčasne je dôležité porovnať viacero malých modelov a posúdiť, či ich doménová adaptácia prináša reálny prínos pre podporu rozhodovania. Výsledok má ukázať, či menšie modely môžu predstavovať efektívnu alternatívu k veľkým všeobecným modelom tam, kde je rozhodujúca dostupnosť, rýchlosť a prevádzková nenáročnosť.


\section{Formulácia úlohy}

Úlohou práce je navrhnúť a implementovať AI asistenta zameraného na oblasť energetiky v Európe, ktorý bude využívať doučené malé jazykové modely a bude schopný poskytovať rýchle, vecné a použiteľné odpovede pre potreby vzdelávania aj podpory rozhodovania. Riešenie má byť realizované tak, aby bolo použiteľné aj pri obmedzených výpočtových zdrojoch.

Na splnenie tejto úlohy sú stanovené tieto čiastkové ciele:

\begin{itemize}
	\item pripraviť a spracovať doménovo orientovaný dataset z energetických dát,
	\item doučiť vybrané malé jazykové modely na pripravených dátach,
	\item navrhnúť jednotný spôsob inferencie a používateľskej interakcie vo forme konverzačného asistenta,
	\item porovnať vybrané modely z pohľadu kvality odpovedí, konzistentnosti a praktickej využiteľnosti,
	\item vyhodnotiť, či doménová adaptácia malých modelov prináša merateľný prínos pre rýchle získanie informácií a podporu rozhodovania.
\end{itemize}

Výsledkom práce má byť funkčný prototyp asistenta a analytické zhodnotenie vhodnosti malých jazykových modelov pre úzko špecializovanú doménu energetiky.
